\clearpage\phantomsection\label{exercise-index-1}
\BookHeading{Percorso degli esercizi guidati}
Quindici esercizi accompagnano lo studio del modulo. Ogni traccia segue le sezioni richieste e rinvia a uno svolgimento completo. Si suggerisce di formulare la previsione e tentare i passaggi prima di consultare la soluzione.
\begin{longtable}{@{}p{13mm}p{75mm}r@{}}\toprule Codice & Argomento & Pagina\\\midrule\endhead
\hyperref[ex-G1.01]{G1.01} & Prevedere due moti dagli stessi dati di posizione & \pageref{ex-G1.01}\\[3pt]
\hyperref[ex-G1.02]{G1.02} & Propagare un'incertezza sulla velocità & \pageref{ex-G1.02}\\[3pt]
\hyperref[ex-G1.03]{G1.03} & Invertire un moto e ritrovare la preparazione & \pageref{ex-G1.03}\\[3pt]
\hyperref[ex-G1.04]{G1.04} & Confrontare ritorno ideale e ritorno perturbato & \pageref{ex-G1.04}\\[3pt]
\hyperref[ex-G1.05]{G1.05} & Sommare due onde con fase assegnata & \pageref{ex-G1.05}\\[3pt]
\hyperref[ex-G1.06]{G1.06} & Dalle ampiezze all'intensità & \pageref{ex-G1.06}\\[3pt]
\hyperref[ex-G1.07]{G1.07} & Localizzare i minimi per una fenditura stretta & \pageref{ex-G1.07}\\[3pt]
\hyperref[ex-G1.08]{G1.08} & Attraversare i regimi della diffrazione & \pageref{ex-G1.08}\\[3pt]
\hyperref[ex-G1.09]{G1.09} & Prevedere la trasmissione di un analizzatore & \pageref{ex-G1.09}\\[3pt]
\hyperref[ex-G1.10]{G1.10} & Seguire una preparazione sulla sfera di Bloch & \pageref{ex-G1.10}\\[3pt]
\hyperref[ex-G1.11]{G1.11} & Preparare, introdurre una fase e ricombinare & \pageref{ex-G1.11}\\[3pt]
\hyperref[ex-G1.12]{G1.12} & Ricavare marginali e correlazioni da una tabella congiunta & \pageref{ex-G1.12}\\[3pt]
\hyperref[ex-G1.13]{G1.13} & Seguire l'effetto di un'interazione con l'ambiente & \pageref{ex-G1.13}\\[3pt]
\hyperref[ex-G1.14]{G1.14} & Stabilire quali dati servano alla previsione successiva & \pageref{ex-G1.14}\\[3pt]
\hyperref[ex-G1.15]{G1.15} & Ricostruire un esperimento a partire dai suoi risultati & \pageref{ex-G1.15}\\[3pt]
\bottomrule\end{longtable}
{\small I primi cinque esercizi offrono una guida esplicita; i successivi cinque consolidano il calcolo; gli ultimi cinque collegano più concetti. Tutti conservano uno svolgimento completo.}\par
